\FloatBarrier
\section{Dynamics of Lin's Model}
\label{sec:lin-results}

Our starting question was whether the five-element network could maintain
an active balanced state after one component was disturbed. A trajectory
that happens to remain at equilibrium does not answer this question: the
important test is whether nearby states return to that equilibrium. We
therefore examined both local stability and the response to initial
perturbations. This section introduces the Lin reference model and the
baseline tests; Section~\ref{sec:modified-results} discusses the modified
model and the broader parameter comparisons.

Throughout, ``balance'' means an equilibrium with equal positive activities
at the five nodes. This is a convenient reference state for a model whose
coefficients are identical across nodes. It is not an assertion that real
organs have equal activities, or that mathematical stability alone defines
health. The experiments below are preliminary comparisons at specified
parameters, rather than calibrated physiological predictions.

\subsection{The reference model and the meaning of balance}
We order the nodes as Wood, Fire, Earth, Metal, and Water, and take indices
modulo five. The predecessor $i-1$ is the mother, $i+1$ is the child,
$i-2$ is the controller, and $i+2$ is the controlled target. In this
ordering, the homogeneous model of Lin and Han~\cite{Lin2020} is
\begin{equation}
 \dot x_i=a x_i-bx_i^2+c x_i x_{i-1}-d x_i x_{i-2}
                 -g x_i x_{i+1}-h x_i x_{i+2}.
 \label{eq:lin-model}
\end{equation}
The terms describe self-growth, self-limitation, nourishment received,
restraint received, a cost associated with the child, and a cost associated
with the controlled target. All five activities are state variables;
``six-term'' describes the equation, not the number of nodes.
For positive coefficients, the positive balanced level is
\begin{equation}
 s_{\mathrm L}=\frac{a}{b+d+g+h-c},
 \qquad b+d+g+h-c>0.
 \label{eq:lin-balance}
\end{equation}
Existence and stability are separate questions. Equation~\eqref{eq:lin-balance}
identifies a possible balance; it does not tell us whether the network
returns there after a disturbance.

\subsection{Can the baseline balance resist a small disturbance?}
We first used the reference choice $(a,b,c,d,g,h)=(1,1,1,1,1,1)$,
which gives $s_{\mathrm L}=1/3$. To test local stability, we evaluated
the Jacobian at $s_{\mathrm L}\mathbf{1}$ and plotted its five eigenvalues.
An eigenvalue with positive real part indicates a direction in which a
small disturbance grows; all negative real parts imply that sufficiently
small disturbances decay. Figure~\ref{fig:baseline-spectra} also includes
the modified-model baseline discussed in Section~\ref{sec:modified-results}
so that the two behaviors can be compared directly.

\begin{figure}[H]
\centering
\includegraphics[width=0.90\linewidth]{figures/figure1.png}
\caption{Local stability at two illustrative baselines. Lin uses
$(a,b,c,d,g,h)=(1,1,1,1,1,1)$, with balanced level $1/3$.
The modified model uses $(a,k,c,g,h)=(2,0.5,1,0.5,0.3)$, with balanced
level approximately $2.670286$. The vertical line separates negative and
positive eigenvalue real parts. The baselines are specified examples,
not a matched test of one changed mechanism.}
\label{fig:baseline-spectra}
\end{figure}

Two Lin eigenvalues lie to the right of the imaginary axis, at approximately
$0.206011\pm0.634038\,\mathrm{i}$. Thus the balanced state exists but is
unstable. The nonzero imaginary parts indicate an oscillatory component in
the growing linearized disturbance, rather than a purely monotone departure.
This gives a concrete expectation for the time-domain experiment:
a generic small displacement from balance should initially produce growing
oscillations. The eigenvalue plot alone does not determine what happens
once the solution has moved far from equilibrium.

\subsection{What happens after one node is disturbed?}
To connect the local calculation with observable trajectories, we started
all five nodes at balance and changed only Wood. The first experiment
increased Wood by $5\%$, testing the response near equilibrium. The second
reduced Wood to $30\%$ of its balanced value, testing a much larger initial
displacement. The other four components started at balance in both cases.
This distinction matters: local stability addresses sufficiently small
perturbations, while the larger disturbance requires a direct trajectory
calculation.

\begin{figure}[H]
\centering
\includegraphics[width=\linewidth]{figures/figure2.png}
\caption{Responses to the same relative initial disturbance at the two
baselines in Figure~\ref{fig:baseline-spectra}. Top: Wood increased by
$5\%$. Bottom: Wood reduced to $30\%$ of balance. Each activity is divided
by its model's balanced level, so the horizontal reference value is one.
All runs extend to $t=120$; the modified-model panels display only
$t\leq12$ to show the transient. Axis ranges differ between panels.}
\label{fig:disturbance}
\end{figure}

In the upper-left panel, the initially small displacement spreads through
the network and grows into a sustained oscillatory pattern over the
observation interval. This agrees with the local instability identified
in Figure~\ref{fig:baseline-spectra}. The components rise and fall at
different times, so the network does not simply move to a new common
activity level. The lower-left panel shows that the larger Wood reduction
also produces persistent departures from the original balance. At $t=120$,
the Euclidean distances to the balanced vector are approximately $0.8776$
and $0.9899$ for the small and large disturbances, respectively.

The conclusion is specific: the all-ones Lin baseline does not restore
these disturbances to its positive balanced state. Oscillations remain
bounded over the displayed interval; their appearance is not evidence of
chaos, permanent extinction, or failure of every Lin parameter choice.
Indeed, keeping the other coefficients at one but increasing self-limitation
to $b=2$ gives a stable positive balance at $s_{\mathrm L}=0.25$.
This observation motivates the parameter sweeps below: a single baseline
cannot represent the behavior of the whole model family.

The trajectory graphs were obtained by integrating the full five-node
equations with DOP853, relative tolerance $10^{-10}$ and absolute tolerance
$10^{-12}$, and sampling 3,001 times over $[0,120]$. Repeating the four
runs with RK45 at the same tolerances changed the sampled activities by
at most $1.1\times10^{-8}$. These checks support the displayed numerical
behavior. The present experiments concern the full network; a comparison
of isolated or externally constrained two-node motifs remains unfinished.

\FloatBarrier
\section{A modified version of the model and its dynamics}
\label{sec:modified-results}

\subsection{Why examine a different interaction law?}
The baseline results raise a modeling question: does changing how a node
receives input and experiences losses alter its ability to return to a
positive balance? In Lin's model, the nourishment term is proportional to
both the mother and the receiving node. Holding the mother fixed, this
input becomes smaller as the receiving node becomes weaker, and vanishes
when that node is empty. One alternative is to make input available even
to a depleted receiver, while reducing it as the receiver's own activity
increases. The modified model used in the following experiments is
\begin{equation}
 \dot x_i=a x_i-kx_i^2+\frac{c x_{i-1}}{1+x_i}
                  -\frac{g x_i}{1+x_{i+2}}-h x_i x_{i-2}.
 \label{eq:modified-model}
\end{equation}
The input $c x_{i-1}/(1+x_i)$ implements this dependence on receiver
activity. The target-dependent loss also changes: holding $x_i$ fixed,
a larger target reduces the magnitude of $g x_i/(1+x_{i+2})$.
Received restraint remains bilinear. These are mathematical modeling
choices whose effects can be explored; the experiments do not establish
them as physiological laws.

There are two important qualifications before comparing the plots.
First, the direct child-loss term has also been omitted, so the modification
changes several mechanisms simultaneously. Second, coefficient names are
not fully shared between the models: received restraint is $d$ in Lin and
$h$ here. Accordingly, the comparison concerns two complete models rather
than one isolated change in functional form. For nonnegative initial
activities and positive coefficients, equation~\eqref{eq:modified-model}
preserves nonnegativity: at $x_i=0$ its derivative is
$c x_{i-1}\geq0$. Thus replenishing a depleted node is compatible with the
chosen nonnegative activity domain.

\subsection{Does the modified baseline return to balance?}
We used $(a,k,c,g,h)=(2,0.5,1,0.5,0.3)$ as an illustrative baseline.
Its positive balanced level is
$s_{\mathrm D}=(3+\sqrt{59})/4\approx2.670286$.
The blue markers in Figure~\ref{fig:baseline-spectra} all lie in the left
half-plane; the largest real part is approximately $-1.153724$.
This predicts decay of sufficiently small disturbances near this balance.

The right-hand panels of Figure~\ref{fig:disturbance} show the corresponding
network response. After Wood is increased, its excess decays while the
other nodes undergo smaller temporary displacements. After Wood is reduced,
it rises toward its original level, while some other nodes temporarily
overshoot or undershoot theirs. Recovery therefore does not mean that
every component moves monotonically or remains unaffected: activity is
redistributed transiently before the network returns to the common balance.

Both tested trajectories return close to the starting equilibrium, with
final distances of approximately $3.3\times10^{-12}$ and
$1.2\times10^{-11}$ at $t=120$. The small-disturbance experiment is
consistent with the eigenvalue prediction. The larger-disturbance experiment
adds evidence that recovery is not restricted to the particular $5\%$
test, but it does not determine the full basin of attraction or establish
global convergence. The different time and vertical scales in the panels
must also be kept in view: the figure demonstrates different outcomes at
the chosen baselines, not a calibrated physiological recovery-speed ratio.

\subsection{Does the contrast persist when parameters change?}
The baseline comparison could depend on the particular coefficients chosen.
To move beyond single examples, we varied nourishment and received restraint
while holding the remaining coefficients fixed. These axes explore the
balance between one positive and one negative contribution to a node's
activity. They do not exhaust the models' parameter spaces.

For each of the $90\times90$ equally spaced grid points, we first found
the positive balanced equilibrium, if one existed. For the modified model,
this is the larger positive root of
$(k+h)s^2+(k+h-a)s-(a+c-g)=0$; Lin uses
Eq.~\eqref{eq:lin-balance}. We then evaluated the Jacobian at that state,
computed all five eigenvalues, and colored the point by the largest real
part, $\alpha$. This produces a local stability map, rather than a map
obtained by guessing the eventual behavior of a simulated trajectory.
Values below $-10^{-9}$ were counted as stable, values above $10^{-9}$ as
unstable, and values within that tolerance as unresolved. The analytic
Jacobians were independently checked by numerical differentiation before
constructing the plots.

\begin{figure}[H]
\centering
\includegraphics[width=\linewidth]{figures/figure3.png}
\caption{Local stability over $c\in[0.3,1.6]$ and received restraint
$r\in[0.1,1.0]$. Lin uses $(a,b,c,d,g,h)=(1.5,0.65,c,r,0.3,0.3)$;
the modified model uses $(a,k,c,g,h)=(1.5,0.65,c,0.3,r)$.
Blue and red indicate negative and positive $\alpha$, respectively;
color limits are clipped. Gray marks absence of a positive balanced
equilibrium. The black line is an interpolated zero contour. Titles use
all grid points as the denominator; inset percentages use only points
admitting positive balance.}
\label{fig:stability-maps}
\end{figure}

The Lin panel contains distinct regimes. The lower-left region admits a
stable positive balance. Moving toward larger nourishment or larger received
restraint within this particular slice crosses into a region where the
balance exists but is unstable. The effect of increasing restraint is
therefore not automatically stabilizing: stronger negative coupling can
coexist with instability of the coupled network. This is an observation
about the specified slice, not a claim that restraint always destabilizes
the model.

The gray corner has a different interpretation. Here
$b+d+g+h-c=1.25+d-c\leq0$, so there is no finite positive balanced level
of the form considered. It must not be counted as another unstable positive
equilibrium. Of the 8,100 Lin grid points, 7,867 admit positive balance
and 1,497 are stable. Hence the two reported percentages are $18.5\%$ of
all grid points and $19.0\%$ of eligible points. Near the gray region,
color clipping also means that the darkest red does not show the exact
magnitude of every positive growth rate.

The modified panel is entirely blue: every grid point in this window admits
a locally stable positive balanced equilibrium. This extends the observed
contrast beyond the single baseline, but only within this parameter slice.
It does not mean that the modified model is stable at every parameter
choice. An earlier, different sweep with $a=2$, $k=0.5$, $g=0.5$,
$c\in[0.05,2.5]$, and $h\in[0.05,2.0]$ had 5,837 stable points out of
8,100 ($72.1\%$). Thus the size of the reported stable region depends on
both the window and the fixed coefficients. The appropriate interpretation
of Figure~\ref{fig:stability-maps} is a broader stable region for the
modified model in this specified experiment, not a universal advantage.

\subsection{How sensitive is the comparison to parameter selection?}
A two-dimensional map still fixes several coefficients. We therefore also
examined randomly selected parameter sets, asking whether the qualitative
contrast persisted under more than one selection rule. The uniform,
ordered, and tiered scenarios sample different parts of parameter space;
they are sensitivity experiments, not estimates of how frequently the
corresponding conditions occur in people.

Each scenario uses 1,200 accepted draws and NumPy's default random-number
generator with seed 3. We paired the models through a shared draw
$(S,A,C,R,T,L)$ for self-growth, self-limitation, nourishment, received
restraint, target-related loss, and child loss. The assignments were
$(a,b,c,d,g,h)_{\mathrm L}=(S,A,C,R,L,T)$ and
$(a,k,c,g,h)_{\mathrm D}=(S,A,C,T,R)$.
Uniform draws are independent on $[0.1,2]$. Ordered draws use the same
ranges conditioned on $S>C>R>T$, leaving $A$ and $L$ unordered.
Tiered draws are independent on $S\in[1,2]$, $A\in[0.3,1]$,
$C\in[0.5,1.5]$, $R\in[0.2,0.8]$, and $T,L\in[0.1,0.5]$.
Using a shared draw makes the assignment explicit, but does not match the
strengths of different interaction laws or supply the child-loss term
missing from the modified model.

\begin{figure}[H]
\centering
\includegraphics[width=\linewidth]{figures/figure4.png}
\caption{Stability frequencies for the three specified parameter-selection
schemes. Left: stable positive balance among draws admitting the selected
positive equilibrium. Right: stable positive balance among all 1,200 draws.
The paired panels distinguish stability conditional on existence from the
combined requirement of existence and stability. Exact counts are listed
in Table~\ref{tab:sampling}.}
\label{fig:sampling}
\end{figure}

The modified model has a higher stability frequency in all three scenarios,
but the size of the difference depends on the sampling rule. For example,
Lin's conditional frequency rises from $22.8\%$ under uniform sampling to
$47.1\%$ under the ordered scenario. Restricting the parameters therefore
changes the apparent comparison substantially; the percentages are not
fixed characteristics of either model.

The two panels also answer different questions. Under uniform sampling,
$94.1\%$ of the modified model's eligible equilibria are stable, but only
$80.7\%$ of all draws admit such a stable positive balance. The difference
arises because 171 of the 1,200 draws have no positive balanced equilibrium.
The left panel asks whether a positive balance is stable once it exists;
the right asks how often the chosen sampling procedure yields a stable
positive balance at all. Reporting only the left panel would conceal this
loss of eligibility. Similarly, the tiered result of $99.0\%$ describes
the selected coefficient ranges, not near-certainty of stability under
arbitrary parameters.

\begin{table}[H]
\centering
\small
\begin{tabular}{llrrr}
\toprule
Scenario & Model & Stable & Positive balance & Total draws\\
\midrule
Uniform & Lin      & 273  & 1195 & 1200\\
        & Modified & 968  & 1029 & 1200\\
Ordered & Lin      & 562  & 1192 & 1200\\
        & Modified & 1103 & 1200 & 1200\\
Tiered  & Lin      & 267  & 1145 & 1200\\
        & Modified & 1188 & 1200 & 1200\\
\bottomrule
\end{tabular}
\caption{Counts underlying Figure~\ref{fig:sampling}. For the modified
model, eligibility and stability refer to the selected upper positive
balanced branch. The two denominators are retained explicitly.}
\label{tab:sampling}
\end{table}

\subsection{What these comparisons establish}
Together, the four graphs address progressively broader questions: whether
balance resists small disturbances, how the full network responds to a
finite displacement, whether the outcome persists across a parameter slice,
and how it changes under different parameter-selection rules. Within the
reported experiments, the modified model admits stable positive balance
more frequently and recovers from both disturbances at its selected
baseline. Lin also admits stable balance, and the modified model is not
stable for every parameter choice.

The observations motivate further comparison; they do not identify one
cause for the difference. The models change several interaction terms,
and the illustrative baselines also use different coefficients. Removing
Lin's child-loss term and then comparing the remaining bilinear model
with the modified model would separate that removal from the changes in
interaction laws. Further comparisons would be needed to distinguish those
laws from one another. Those controlled experiments are separate from the
original Lin comparison reported here. The present figures should therefore
be read as preliminary evidence about the behavior of two complete models
under stated assumptions. Analysis of the associated two-node motifs
remains to be completed.
\FloatBarrier

